\documentclass[11pt]{article}
\usepackage{amsmath, amssymb, graphicx, hyperref}
\usepackage{enumitem}
\setlist{nosep}
\usepackage[margin=1in]{geometry}

\title{In-Class Problem Set: Uncertainty and Publication-Quality Figures with World Development Indicators (R + GitHub \emph{or} Canvas)}
\author{}
\date{}

\begin{document}
\maketitle

\noindent \textbf{Goal.} Practice communicating uncertainty and applying the publishing-quality checklist using real political and economic data. You will build error bar plots, uncertainty ribbons, and redesign a rough figure into a polished version. You may submit via \textbf{GitHub or Canvas}.

\medskip
\noindent \textbf{Dataset.} This problem set uses two sources:
\begin{itemize}
  \item The \texttt{gapminder} package (country-level panel data on life expectancy, GDP per capita, and population).
  \item Simulated survey data you will create in R (following the structure from lecture).
\end{itemize}

\medskip
\noindent \textbf{Key variables (gapminder).}
\begin{itemize}
  \item \texttt{country}, \texttt{continent}, \texttt{year}
  \item \texttt{lifeExp}: Life expectancy at birth
  \item \texttt{pop}: Population
  \item \texttt{gdpPercap}: GDP per capita (inflation-adjusted)
\end{itemize}

\medskip
\noindent \textbf{What to submit (GitHub or Canvas).}
\begin{itemize}
  \item \texttt{scripts/lab.R}
  \item \texttt{outputs/writeup.md}
  \item Required figures in \texttt{figures/}
\end{itemize}

\medskip
\noindent \textbf{Rules.}
\begin{itemize}
  \item Work inside an \textbf{R Project}.
  \item Use a \textbf{sequential, hard-coded workflow} (no user-defined functions).
  \item Save all plots using \texttt{ggsave()}.
  \item If submitting via GitHub, Git commands must be run in the \textbf{Terminal tab}.
\end{itemize}

\section*{Questions}

\begin{enumerate}

\item \textbf{Load data and create simulated survey estimates (proof required).}

\begin{verbatim}
# load required packages
library(__________)
library(__________)
library(__________)

# gapminder data
library(gapminder)
gap <- gapminder

# create simulated survey data
set.seed(42)
df_survey <- data.frame(
  country = c("United States", "United Kingdom", "Germany",
              "France", "Japan", "Brazil", "India", "Nigeria"),
  dem_support = c(72, 68, 65, 61, 54, 48, 42, 38),
  n = c(1200, 1000, 950, 1100, 800, 600, 750, 500)
)
df_survey$se <- sqrt(df_survey$dem_support *
  (100 - df_survey$dem_support) / df_survey$n)
df_survey$lower <- df_survey$dem_support - 1.96 * df_survey$se
df_survey$upper <- df_survey$dem_support + 1.96 * df_survey$se

dim(gap)
head(df_survey)
\end{verbatim}

Include in your write-up:
\begin{itemize}
  \item number of rows and columns in \texttt{gapminder},
  \item the range of years,
  \item confirmation that \texttt{df\_survey} has 8 rows with \texttt{estimate}, \texttt{se}, \texttt{lower}, \texttt{upper}.
\end{itemize}

\item \textbf{Error bar plot: support for democracy.}

Create a dot plot of \texttt{dem\_support} by \texttt{country} with 95\% confidence intervals.

You must:
\begin{itemize}
  \item Reorder countries by estimate.
  \item Add error bars using \texttt{geom\_errorbarh()}.
  \item Add a reference line at 50\% (majority threshold).
  \item Use \texttt{theme\_classic()} and clear labels.
\end{itemize}

\textbf{Pseudo-code structure:}

\begin{verbatim}
df_survey$country <- reorder(df_survey$country, df_survey$dem_support)

ggplot(df_survey, aes(x = ________, y = ________)) +
  geom_vline(xintercept = 50, linetype = "dashed") +
  geom_errorbarh(aes(xmin = ________, xmax = ________),
                 height = 0.25) +
  geom_point(size = 3) +
  theme_classic() +
  labs(x = "________", y = NULL, title = "________")
\end{verbatim}

Save as:
\begin{itemize}
  \item \texttt{figures/democracy\_support\_ci.png}
\end{itemize}

\item \textbf{Time series with uncertainty ribbon.}

Using \texttt{gapminder}, compute the mean and standard error of \texttt{lifeExp} by \texttt{continent} and \texttt{year}. Then create a time series with uncertainty ribbons.

\textbf{Pseudo-code structure:}

\begin{verbatim}
gap_summary <- gap %>%
  group_by(________, ________) %>%
  summarise(
    mean_le = mean(________),
    se_le = sd(________) / sqrt(n()),
    lower = mean_le - 1.96 * se_le,
    upper = mean_le + 1.96 * se_le,
    .groups = "drop"
  )

ggplot(gap_summary, aes(x = year, color = ________, fill = ________)) +
  geom_ribbon(aes(ymin = lower, ymax = upper), alpha = 0.15, color = NA) +
  geom_line(aes(y = mean_le), linewidth = 0.7) +
  theme_classic() +
  theme(legend.position = "bottom") +
  labs(x = NULL, y = "________", color = NULL, fill = NULL,
       title = "________")
\end{verbatim}

Save as:
\begin{itemize}
  \item \texttt{figures/life\_exp\_ribbon.png}
\end{itemize}

\item \textbf{Redesign exercise: before and after.}

Create a deliberately rough ``before'' version of any figure from this lab (no title, default theme, no units, no caption). Then create a polished ``after'' version applying the full publishing-quality checklist:
\begin{itemize}
  \item Descriptive title and subtitle
  \item Axis labels with units
  \item Caption with data source
  \item Reference lines if appropriate
  \item Color for emphasis (not decoration)
  \item Readable text size (\texttt{base\_size} argument)
  \item Clean theme
\end{itemize}

Save as:
\begin{itemize}
  \item \texttt{figures/redesign\_before.png}
  \item \texttt{figures/redesign\_after.png}
\end{itemize}

\item \textbf{Regression uncertainty: GDP and life expectancy.}

Using \texttt{gapminder} data from 2007 only, create a scatterplot of \texttt{gdpPercap} (x-axis, log scale) vs \texttt{lifeExp} (y-axis) with a linear smoother and 95\% confidence ribbon.

\begin{verbatim}
gap_2007 <- gap %>% filter(year == 2007)

ggplot(gap_2007, aes(x = ________, y = ________)) +
  geom_point(alpha = 0.5, aes(size = pop)) +
  geom_smooth(method = "lm", se = TRUE) +
  scale_x_log10(labels = dollar) +
  scale_size_continuous(labels = comma, guide = "none") +
  theme_classic() +
  labs(x = "________", y = "________", title = "________")
\end{verbatim}

Save as:
\begin{itemize}
  \item \texttt{figures/gdp\_lifeexp\_lm\_se.png}
\end{itemize}

\item \textbf{Interpretation (write-up required).}

Write 12--16 sentences addressing:
\begin{itemize}
  \item In the democracy support plot, which countries have overlapping confidence intervals? What does that mean for ranking them?
  \item How does the uncertainty ribbon on the life expectancy time series change your reading of continent-level trends?
  \item In your before/after redesign, name three specific changes you made and why each improved the figure.
  \item In the GDP--life expectancy regression, what does the confidence ribbon around the fitted line represent in plain language?
  \item Name one figure from this lab where uncertainty was especially important to show and explain why.
\end{itemize}

\item \textbf{Submit your work (GitHub or Canvas; proof required).}

\begin{itemize}
  \item \textbf{GitHub option:}
\begin{verbatim}
git status
git add .
git commit -m "Uncertainty lab: CIs + ribbons + redesign"
git push
\end{verbatim}

  \item \textbf{Canvas option:}
  Upload required files to Canvas.
\end{itemize}

Include proof in write-up:
\begin{itemize}
  \item GitHub: output of \texttt{git status} and \texttt{git log -1}
  \item Canvas: confirmation note + list of uploaded files
\end{itemize}

\end{enumerate}

\section*{Optional Extension: Shiny App}

Build a minimal Shiny app that lets the user choose a continent from a dropdown and displays the life expectancy time series (with uncertainty ribbon) for that continent only.

\begin{verbatim}
library(shiny)

ui <- fluidPage(
  selectInput("continent", "Choose continent:",
              choices = unique(gap_summary$continent)),
  plotOutput("lePlot")
)

server <- function(input, output) {
  output$lePlot <- renderPlot({
    df_filtered <- gap_summary %>% filter(continent == input$continent)
    ggplot(df_filtered, aes(x = year, y = mean_le)) +
      geom_ribbon(aes(ymin = lower, ymax = upper), fill = "grey70", alpha = 0.4) +
      geom_line(linewidth = 0.8) +
      theme_classic() +
      labs(x = NULL, y = "Life expectancy (years)",
           title = paste(input$continent, ": life expectancy over time"))
  })
}

shinyApp(ui, server)
\end{verbatim}

\section*{Checklist}
\begin{itemize}
  \item Script runs top-to-bottom
  \item Required figures saved
  \item Write-up complete
  \item Submitted via GitHub or Canvas
\end{itemize}

\end{document}
